% ============================================================
% CHAPTER 2
% ============================================================

\chapter{Collapse as Derived Truncation}
\label{chap:collapse-truncation}

\section{Motivation}

The second fundamental operation in the SpherePop calculus is
\emph{collapse}: an abstraction mechanism which replaces a geometric
configuration by a simpler configuration in which certain internal
degrees of freedom have been removed. Intuitively, collapse forgets
fine--scale detail while preserving large--scale structure and boundary
information. In the lamphrodynamic setting, collapse corresponds to an
entropy--preserving coarse--graining of lamphrodynamic fields.

Heuristically, collapse should behave like an $\infty$--categorical
truncation, removing higher homotopy levels while retaining essential
boundary invariants. The central conjecture of this chapter asserts
that collapse \emph{is} such a truncation in the derived configuration
stack $\dConfig$.

\section{Collapse as a Functor}

Let $\Collapse$ denote the SpherePop collapse operation acting on the
(enriched) configuration category. Formally, collapse is a functor
\[
  \Collapse \colon \dConfig \longrightarrow \dConfig,
\]
equipped with a natural transformation $\Collapse \Rightarrow \id$
generated by the embedding of collapsed configurations into their
refinements. Intuitively, $\Collapse(X)$ identifies lamphrodynamically
degenerate internal structure in $X$ and retracts it to a simpler
configuration which preserves entropic boundary conditions.

\begin{definition}
A collapse is \emph{entropy--admissible} if it preserves the
restriction of $\SField$ to $\partial X$, i.e.\ if
$\SField|_{\partial X}$ pulls back along the collapse map. We require
entropy--admissible collapse maps throughout.
\end{definition}

\begin{remark}
Entropy--admissibility is the minimal requirement for collapse to
represent a physically meaningful coarse--graining in lamphrodynamic
evolution. Stronger admissibility conditions may be imposed to preserve
additional lamphrodynamic invariants.
\end{remark}

\section{Connection to $\infty$--Categorical Truncation}

Let $\tau_{\leq k}$ denote the $k$--truncation functor on an
$\infty$--category. The guiding idea is that collapse removes
higher--order lamphrodynamic fluctuations (as measured by derived
obstruction theory), corresponding to truncation at some effective
level $k$.

\begin{conjecture}[Collapse as Derived Truncation]
\label{conj:collapse-truncation}
There exists an integer $k \geq 0$, depending on the lamphrodynamic
regularity of $X$, such that
\[
  \Collapse(X) \simeq \tau_{\leq k}(X)
\]
in the derived configuration stack $\dConfig$.
\end{conjecture}

\begin{remark}
The value of $k$ depends on lamphrodynamic smoothness and boundary
entropy. In particular, configurations with higher lamphrodynamic
complexity require deeper truncation levels to preserve admissible
boundary data.
\end{remark}

\section{Evidence for the Conjecture}

Evidence for \cref{conj:collapse-truncation} arises from three sources:

\begin{enumerate}
\item \emph{Derived obstruction theory.} Internal lamphrodynamic modes
correspond to higher obstruction classes in the tangent complex of
$\dConfig$. Collapse eliminates modes associated to these obstructions,
suggesting truncation of higher homotopy groups.

\item \emph{Entropy monotonicity.} Collapse preserves boundary entropy
while reducing internal entropy, consistent with truncation of
non--boundary degrees of freedom.

\item \emph{Computational abstraction.} In SpherePop--based computation,
collapse produces coarser representations which retain semantic boundary
structure but forget internal detail---again consistent with removing
higher homotopy levels.
\end{enumerate}

\section{Partial Results}

Although we cannot prove \cref{conj:collapse-truncation} in full
generality, we can establish partial results in special cases.

\begin{proposition}
If $X$ is lamphrodynamically regular and its entropy density $\SField$
is constant on connected components, then $\Collapse(X)$ is equivalent
to the $0$--truncation of $X$ in $\dConfig$.
\end{proposition}

\begin{proof}[Sketch]
Under these assumptions, internal lamphrodynamic variations vanish, so
higher obstruction classes in the tangent complex are trivial. Hence
the derived structure is concentrated in degree $0$, and collapse
reduces to a $0$--truncation.
\end{proof}

\section{Boundary Invariants}

A key requirement for collapse is the preservation of lamphrodynamic
boundary conditions. These provide invariants under collapse, analogous
to homotopy invariants under truncation.

\begin{definition}
A \emph{boundary invariant} of $X$ is any invariant of the restricted
lamphrodynamic field $(\PhiField,\vField,\SField)|_{\partial X}$.
\end{definition}

\begin{remark}
Boundary invariants classify equivalence classes of collapsed
configurations, and hence determine semantic or physical information
preserved under computational abstraction.
\end{remark}

\section{Open Problems}

\begin{enumerate}
\item Determine precise conditions under which $\Collapse$ implements
exact $k$--truncation.
\item Classify boundary invariants and their lamphrodynamic dependence.
\item Establish computational complexity bounds for collapse viewed as
entropy--admissible truncation.
\end{enumerate}

These questions link geometric computation to derived category theory
and lamphrodynamic physics, and will be revisited in
Part~III of this volume.
\newpage

